%% CUMCM-Master LaTeX 模板
%% 适用于全国大学生数学建模竞赛
%% 编译方式：xelatex main.tex（建议运行两次以更新交叉引用）

\documentclass[12pt, a4paper]{article}

%% ======================== 宏包导入 ========================
\usepackage{ctex}                    % 中文支持
\usepackage[margin=2.5cm]{geometry}  % 页边距
\usepackage{amsmath, amssymb, amsthm}% 数学公式
\usepackage{graphicx}                % 插图
\usepackage{subcaption}              % 子图
\usepackage{booktabs}                % 三线表
\usepackage{longtable}               % 长表格
\usepackage{array}                   % 表格增强
\usepackage{multirow}                % 跨行单元格
\usepackage{float}                   % 浮动体控制 [H]
\usepackage{hyperref}                % 超链接
\usepackage{enumerate}               % 列表
\usepackage{enumitem}                % 列表定制
\usepackage{listings}                % 代码高亮
\usepackage{xcolor}                  % 颜色
\usepackage{tikz}                    % 流程图
\usepackage{pgfplots}                % 数据可视化
\usepackage{caption}                 % 标题定制
\usepackage{fancyhdr}                % 页眉页脚
\usepackage{setspace}                % 行距
% natbib 与本模板的 plain \bibitem 参考文献不兼容（natbib 要求 author-year
% 格式的 \citep/\citet，直接编译会报 "Bibliography not compatible with
% author-year citations" 致命错误）；模板正文全程只用标准 \cite，不需要
% natbib，故不加载（AutoMCM_SOP §17 格式合规要求）——这是修复模板本身，
% 不是"生成论文时再手动删掉"的临时补丁。
\usepackage{appendix}                % 附录

\pgfplotsset{compat=1.18}
\usetikzlibrary{shapes.geometric, arrows.meta, positioning, calc}

%% ======================== 代码样式 ========================
\definecolor{codegray}{rgb}{0.5,0.5,0.5}
\definecolor{codepurple}{rgb}{0.58,0,0.82}
\definecolor{codeblue}{rgb}{0.13,0.29,0.53}
\definecolor{backcolour}{rgb}{0.97,0.97,0.97}

\lstdefinestyle{pythonstyle}{
    backgroundcolor=\color{backcolour},
    commentstyle=\color{codegray}\itshape,
    keywordstyle=\color{codeblue}\bfseries,
    stringstyle=\color{codepurple},
    basicstyle=\ttfamily\footnotesize,
    breakatwhitespace=false,
    breaklines=true,
    captionpos=b,
    keepspaces=true,
    numbers=left,
    numbersep=5pt,
    numberstyle=\tiny\color{codegray},
    showspaces=false,
    showstringspaces=false,
    showtabs=false,
    tabsize=4,
    frame=single,
    rulecolor=\color{black!20},
    language=Python
}

%% ======================== 页面样式 ========================
\pagestyle{fancy}
\setlength{\headheight}{14pt}
\fancyhf{}
\fancyhead[L]{\small 全国大学生数学建模竞赛}
\fancyhead[R]{\small \leftmark}
\fancyfoot[C]{\thepage}
\renewcommand{\headrulewidth}{0.4pt}

\captionsetup{font=small, labelfont=bf}
\setstretch{1.5}

%% ======================== 自定义命令 ========================
\newtheorem{assumption}{假设}[section]
\newtheorem{definition}{定义}[section]
\newtheorem{theorem}{定理}[section]

\newcommand{\prob}[1]{\mathbf{P}\left(#1\right)}
\newcommand{\E}[1]{\mathbb{E}\left[#1\right]}
\newcommand{\norm}[1]{\left\|#1\right\|}

%% ======================== 文档信息 ========================
\title{
    \vspace{-1cm}
    {\Large \textbf{CUMCM 20XX 竞赛题目}} \\[0.5em]
    {\large 全国大学生数学建模竞赛参赛论文}
}
\author{}
\date{}

%% ======================== 格式合规提醒（AutoMCM_SOP.md §17，不渲染进论文）========
%% 依据《全国大学生数学建模竞赛论文格式规范（2026年修订稿）》：
%%   - 正文（不含附录）不超过 30 页，不要目录
%%   - 摘要专用页开始编页码（阿拉伯数字，页脚居中），本模板已配置
%%   - 电子版提交件第一页必须是摘要专用页——承诺书、编号专用页不放进电子版，
%%     下方已注释掉这两页，只在需要打印纸质版留存时取消注释
%%   - 论文摘要页/正文/附录任何地方不能出现参赛者身份、学校、赛区信息
%%     （quality_gate.py anon-check 会做启发式扫描，但不能完全替代人工复查）
%% ======================================================================

%% ======================== 正文开始 ========================
\begin{document}

%% ---- 纸质版留存专用：承诺书 + 编号专用页 ----------------------------------
%% 电子版提交前必须保持注释状态（规范第十条），仅打印纸质版留存时取消注释：
% \begin{titlepage}
% \thispagestyle{empty}
% \begin{center}
% \vspace*{2cm}
% {\Large 全国大学生数学建模竞赛承诺书}\\[2em]
% （此处填写官方承诺书标准文本，各赛区可能略有差异，以当年官方模板为准）
% \end{center}
% \newpage
% \thispagestyle{empty}
% \begin{center}
% \vspace*{2cm}
% {\Large 编号专用页}\\[2em]
% （此页由承办方填写赛区、题号、编号，参赛队不要自行填写）
% \end{center}
% \end{titlepage}
%% ---------------------------------------------------------------------------

\pagenumbering{arabic}
\setcounter{page}{1}

\maketitle
\thispagestyle{plain}

%% ---------------------- 摘要（电子版第 1 页，从这页开始编页码）--------------
\begin{abstract}
\noindent
（正文结果和证据账本稳定后再撰写摘要。概括任务、模型单元、求解方法、主要结论与验证；所有数字和比较必须能回指已验证证据。原则上不超过一页，不需要翻译成英文。）

\vspace{1em}
\noindent\textbf{关键词：}数学规划；%
% 在此补充关键词，用分号分隔，3~5个
\end{abstract}

\newpage
%% 官方规范明确要求"不要目录"——\tableofcontents 不得出现在正文里

%% ==================== 1. 问题重述 ====================
\section{问题重述}

（在此精炼赛题核心内容。需提炼题目的研究对象、给定条件和核心目标，不超过半页。避免直接复制原题，需用自己的语言重新表述，体现对题意的精准把握。）

%% ==================== 2. 问题背景与需要解决的问题 ====================
\section{问题背景与需要解决的问题}

\subsection{问题背景}

（阐述题目所处的宏观背景，说明该问题的现实意义与研究价值，1~2段。）

\subsection{需要解决的问题}

\begin{enumerate}[label=\textbf{问题\arabic*：}, leftmargin=*]
    \item （问题一描述）
    \item （问题二描述）
    \item （问题三描述，按实际题目增删）
\end{enumerate}

%% ==================== 问题分析与模型单元 ====================
\section{问题分析}

\subsection{小问依赖与模型单元}

（依据 \texttt{problem\_graph} 说明小问之间的共享状态、输入输出接口和依赖，并解释模型单元的合并或拆分理由。）

% 流程图示例（根据实际内容修改）
\begin{figure}[H]
\centering
\begin{tikzpicture}[
    box/.style={rectangle, rounded corners, draw=black, fill=blue!10,
                text width=2.2cm, align=center, minimum height=0.8cm},
    arrow/.style={-Stealth, thick}
]
    \node[box] (input)  {读取数据};
    \node[box, right=0.6cm of input]  (preprocess) {数据预处理};
    \node[box, right=0.6cm of preprocess] (model) {建立模型};
    \node[box, right=0.6cm of model]  (solve) {求解};
    \node[box, right=0.6cm of solve]  (output) {输出结果};

    \draw[arrow] (input) -- (preprocess);
    \draw[arrow] (preprocess) -- (model);
    \draw[arrow] (model) -- (solve);
    \draw[arrow] (solve) -- (output);
\end{tikzpicture}
\caption{模型单元与数据流}
\label{fig:flow1}
\end{figure}

图~\ref{fig:flow1} 应展示模型单元之间真实存在的数据依赖；若流程可以用一两句话说清，则删除该图。

\subsection{机制、约束与求解路线}

（从数据/现象生成机制出发，说明约束、最小基线、基线失效条件和求解器选择；不得仅以算法名称或复杂度作为理由。）

%% ==================== 4. 模型假设 ====================
\section{模型假设}

为使模型具有可操作性，结合题目的实际背景，本文提出以下假设：

\begin{enumerate}[label=\textbf{假设\arabic*：}, leftmargin=*]
    \item \label{assump:h1}（假设内容）。\\
          \textit{依据：}（物理、数据或数学依据。）\\
          \textit{影响与边界：}（该假设简化了什么，何时可能失效。）

    % 按实际建模需要增删；每条假设必须在模型、验证或局限分析中被引用
\end{enumerate}

%% ==================== 5. 符号说明 ====================
\section{符号说明}

本文所用主要符号及其含义如表~\ref{tab:symbols}~所示。

\begin{table}[H]
\centering
\caption{主要符号说明}
\label{tab:symbols}
\begin{tabular}{cll}
\toprule
\textbf{符号} & \textbf{含义} & \textbf{单位/量纲} \\
\midrule
$x_i$  & （变量含义）     & （单位） \\
$y_j$  & （变量含义）     & （单位） \\
$\alpha$ & （参数含义）   & 无量纲 \\
$T$    & 时间周期         & s \\
$n$    & 样本数量         & 个 \\
% 根据实际情况增删行
\bottomrule
\end{tabular}
\end{table}

%% ==================== 模型单元的建立与求解 ====================
\section{模型的建立与求解}

\subsection{模型单元 M1：（名称）}

\subsubsection{机理分析}

（阐述该单元覆盖的小问、物理/数据机理、输入输出接口及与其他单元的关系，并明确说明所用假设，例如假设~\ref{assump:h1}。）

\subsubsection{模型建立}

（写出数学模型，包括目标函数和约束条件。）

\begin{equation}
    \min_{x} \quad f(x) = \sum_{i=1}^{n} w_i x_i
    \label{eq:obj1}
\end{equation}

\begin{equation}
    \text{s.t.} \quad
    \begin{cases}
        g_i(x) \leq 0, & i = 1, \ldots, m \\
        h_j(x) = 0,    & j = 1, \ldots, p \\
        x_i \geq 0,    & i = 1, \ldots, n
    \end{cases}
    \label{eq:constraint1}
\end{equation}

\subsubsection{基线、失效与求解算法}

（先给最小可解释基线及其失效条件，再描述求解器、参数、相对基线的可验证收益和回退方案。）

\subsubsection{求解结果}

（展示求解结果，用表格和图形呈现。）

\begin{table}[H]
\centering
\caption{问题一求解结果}
\label{tab:result1}
\begin{tabular}{cccc}
\toprule
\textbf{指标} & \textbf{数值} & \textbf{单位} & \textbf{说明} \\
\midrule
（指标名）  & （数值）  & （单位）  & （说明） \\
\bottomrule
\end{tabular}
\end{table}

\begin{figure}[H]
    \centering
    \IfFileExists{images/fig01_result1.png}{%
        \includegraphics[width=0.75\textwidth]{images/fig01_result1.png}%
    }{\fbox{\parbox[c][4cm][c]{0.7\textwidth}{\centering 运行模型后替换为实际结果图}}}
    \caption{（图题）}
    \label{fig:result1}
\end{figure}

（结合图~\ref{fig:result1} 和相应数值表分析结果，说明实际意义及对后续小问的作用。）

\subsection{模型单元 M2：（名称，按 \texttt{problem\_graph} 增删）}

（说明它复用哪些状态或接口、改变了哪些机制/约束，以及回答哪些小问。）

%% ==================== 7. 模型的分析与检验 ====================
\section{模型的分析与检验}

\subsection{误差、敏感性与稳健性分析（按模型需要选择）}

（围绕主要失效源选择检查。扰动范围应来自测量误差、参数区间、数值容差或有依据的情景边界；不为填表而使用固定比例。）

\begin{table}[H]
\centering
\caption{灵敏度分析结果}
\label{tab:sensitivity}
\begin{tabular}{llll}
\toprule
\textbf{对象} & \textbf{范围/网格} & \textbf{依据} & \textbf{关键结果} \\
\midrule
（参数/步长/情景） & （实际范围） & （误差或尺度来源） & （结果） \\
\bottomrule
\end{tabular}
\end{table}

\subsection{模型单元验证}

（按模型类型报告初边值/约束残差、收敛、简化基准、交叉方法或不变量；每项写明阈值依据和证据位置。）

%% ==================== 8. 模型的评价、改进与推广 ====================
\section{模型的评价、改进与推广}

\subsection{模型的有效性与优点}

\begin{enumerate}
    \item （结合具体机制、验证结果或计算代价说明；按实际需要增删。）
\end{enumerate}

\subsection{模型局限性}

\begin{enumerate}
    \item （点名具体假设、数据限制或算法失效条件；按实际需要增删。）
\end{enumerate}

\subsection{改进方向与推广}

（仅写有明确基线、失效点和验证方案的改进；若无可靠证据，可只说明适用边界。）

%% ==================== 9. 参考文献 ====================
\section*{参考文献}
\begin{thebibliography}{99}

\bibitem{ref1}
（作者姓名）.（文献标题）[J/M/C].（期刊/出版社）,（年份）,（卷号）(（期号）):（页码）.

\bibitem{ref2}
Author A, Author B. Title of the paper[J]. \textit{Journal Name}, Year, Vol(Issue): Pages. DOI: xxx.

% 只列实际使用并完成真实性核验的文献，不按数量凑条目

\end{thebibliography}

%% ==================== 附录 ====================
\appendix
\section{支撑材料文件清单}

（逐项列出相对路径、用途及其生成/验证的结果文件。清单必须覆盖全部可运行源程序和题目要求的支撑材料。）

\begin{longtable}{p{0.42\textwidth}p{0.48\textwidth}}
\toprule
\textbf{相对路径} & \textbf{用途} \\
\midrule
\texttt{src/models/...} & （模型单元、覆盖小问及输出） \\
\texttt{src/verifications/...} & （验证对象及报告） \\
\bottomrule
\end{longtable}

\section{完整可运行源程序}

% 最终稿必须逐个列出 src/ 下所有可运行源文件；不得设置 firstline/lastline。
% 下面命令在对应文件存在后取消注释并按实际路径增删。

\subsection{数据预处理（\texttt{01\_data\_eda.py}）}
% \lstinputlisting[style=pythonstyle, caption={数据探索与预处理}]{../src/models/01_data_eda.py}

\subsection{模型单元（按实际文件列出）}
% \lstinputlisting[style=pythonstyle, caption={模型单元 M1}]{../src/models/model_unit1.py}

\subsection{验证程序（按实际文件列出）}
% \lstinputlisting[style=pythonstyle, caption={模型单元 M1 验证}]{../src/verifications/verify_model_unit1.py}

\end{document}
